\documentclass[10pt,a4paper]{amsart}
\usepackage[margin=19mm]{geometry}
\usepackage{amsmath,amssymb,mathtools}
\usepackage[scr=rsfs,cal=euler]{mathalfa}
\usepackage{xfrac}
\usepackage[unicode,hidelinks,bookmarks=false]{hyperref}
\newcommand{\ie}{i.e.\ }
\newcommand{\cf}{cf.\ }
\newcommand{\resp}{respectively\ }
\newcommand{\Subsection}{\subsection}
\providecommand{\nobreakdash}{\nobreak\hbox{-}}
\setlength{\emergencystretch}{2em}
\multlinegap0pt
\usepackage{amssymb}
\usepackage{xcolor}
\usepackage{enumerate}
\usepackage{tikz}
\newtheorem{theorem}{Theorem}
\newtheorem{lemma}{Lemma}
\newcommand{\BF}{{\boldsymbol{F}}}
\newcommand{\Be}{{\boldsymbol{e}}}
\newcommand{\Bu}{{\boldsymbol{u}}}
\newcommand{\Bv}{{\boldsymbol{v}}}
\newcommand{\R}{\mathbb{R}}
\newcommand{\curl}{\mathrm{curl}}
\newcommand{\mcC}{\mathcal{C}}
\newcommand{\n}{\rho}
\newcommand{\p}{\partial}
\newcommand{\rmd}{\mathrm{d}}
\newcommand{\te}{\theta}
\newcommand{\vp}{\varphi}
\title{Forced self-similar solutions to the stationary Navier--Stokes equations in a half-space}
\author{Yun Wang, Chunjing Xie, Shaoheng Zhang}
\date{Source: arXiv:2608.25432v1 (CC BY 4.0)}
\begin{document}
\maketitle

\noindent\emph{Source and licence.} Forced self-similar solutions to the stationary Navier--Stokes equations in a half-space. Version arXiv:2608.25432v1 (CC BY 4.0), distributed by arXiv under the Creative Commons Attribution 4.0 International licence, http://creativecommons.org/licenses/by/4.0/, as recorded in the arXiv licence metadata for this exact version and shown on the abstract page https://arxiv.org/abs/2608.25432v1. Full licence notice: \texttt{licenses/arxiv-2608.25432-CC-BY-4.0.txt}.\par\smallskip

\noindent\textit{Editorial scope.} Counted unit: the article's theorem thm:A (source0.tex lines 313--327). The article's complete original proof of it is delivered with it (source0.tex lines 831--900, one \texttt{proof} environment of 70 lines, which the article places away from the statement). No mathematics is written, repaired or re-derived: the statement and the proof are the article's own lines, in the article's own notation, and no retained line is substituted or re-flowed.  Retained uncounted as the source-local context the proof needs: lines 211--225; lines 558--571; lines 594--603; lines 623--631; lines 633--635; lines 682--690; lines 706--713; lines 742--759; lines 1505--1505.  Nothing was omitted from any retained span, so the package carries no omission record.  No retained line contains a citation, so the wrapper carries no bibliography and no bibliography entry.  No retained line contains a figure environment, an image-inclusion command or drawing code, so no image is delivered and the package declares no asset dependency.  Editorial formatting only, in addition: an AMS \texttt{amsart} preamble that restates the article's own declarations and macros used by the retained lines, namely the environments \texttt{theorem}, \texttt{lemma} on separate counters, together with the standard packages those lines need.  The retained environments print in this document as Theorem 1, Lemma 1 in order of appearance, whereas the article prints the same items with the numbering of its own full document.  The complete native source of the article as distributed in the arXiv e-print for this exact version is bundled unchanged as \texttt{sources/208579/source0.tex}.
\begin{equation}\label{eqn:NSEhfD}
	\left\{
	\begin{aligned}
		-\Delta \Bu+(\Bu\cdot\nabla)\Bu+\nabla p
		&=\BF
		&&\text{in }\mathbb R^3_+,\\
		\operatorname{div}\Bu
		&=0
		&&\text{in }\mathbb R^3_+,\\
		\Bu
		&=0
		&&\text{on }\partial\mathbb R^3_+\setminus\{0\}.
	\end{aligned}
	\right.
\end{equation}

\begin{equation}\label{eqn:sphericalcurl}
\begin{aligned}
\curl \, \Bv=
\frac{1}{\rho\sin\vp}
\left(
\frac{\p(v^{\te}\sin\vp)}{\p\vp}-\frac{\p v^{\vp}}{\p\te}
\right)
\Be_{\rho}
+\frac{1}{\rho}
\left(\frac{1}{\sin\vp}\frac{\p v^{\rho}}{\p\te}-\frac{\p(\rho v^{\te})}{\p\rho}\right)\Be_{\vp}
+\frac{1}{\rho}
\left(\frac{\p(\rho v^{\vp})}{\p\rho}-\frac{\p v^{\rho}}{\p\vp}\right)\Be_{\te},  
\end{aligned}
\end{equation}

\begin{equation}\label{eqn:halfD1}
\left\{
\begin{aligned}
   f^{\prime\prime} + f^{\prime} \cot \varphi &= gf^{\prime} -(f^2 + g^2 + h^2) - 2P - X, \\
   f^{\prime} &= gg^{\prime} - h^2 \cot \varphi + P^{\prime} - Y, \\ 
   (h^\prime + h\cot \varphi)^{\prime} &= g(h^\prime + h\cot \varphi) - Z,\\
   f+ g^\prime + g\cot \varphi &=0,
\end{aligned}
\right.
\end{equation}

\begin{equation}\label{eqn:halfD2}
\left\{ 
\begin{aligned}
    f^{\prime \prime} + f^\prime \cot \varphi &= f^\prime g - f^2 - ( g^2+2P) -X,\\
    f^\prime &= g g^\prime + P^\prime - Y, \\
    f+ g^\prime + g\cot \varphi &=0. 
\end{aligned}
\right.
\end{equation}

\begin{align}\label{eqn:eqnforP}
\frac12 g^2 + P = f + \int_0^{\varphi} Y(\tau)\,\mathrm{d}\tau + \mcC
\end{align}

\begin{equation}\label{eqn:constC}
\begin{aligned}
\int_0^1 J(s) \, \mathrm{d} s 
&= \int_0^{\frac{\pi}{2}} \biggl(\int_0^{\varphi} (G(\tau) +2 \mcC) \sin\tau \, \mathrm{d}\tau \biggr) \sin\varphi \, \mathrm{d}\varphi\\
&= \int_0^{\frac{\pi}{2}} G(\tau) \sin\tau \biggl(\int_{\tau}^{\frac{\pi}{2}}
\sin\varphi \, \mathrm{d}\varphi \biggr) \, \mathrm{d}\tau + \mcC \int_0^{\frac{\pi}{2}} 2\sin\tau(1-\cos\tau)\,\mathrm{d}\tau\\
&= \int_0^{\frac{\pi}{2}} G(\tau) \sin\tau \cos\tau \, \mathrm{d}\tau + \mcC = 0.
\end{aligned}
\end{equation}

\begin{equation}\label{eqn:Riccati}
\left\{
\begin{aligned}
\frac{\mathrm{d} M}{\mathrm{d} t} + \frac12 M^2 &= N \quad \text{in } (0,1), \\
M(0)&=0.
\end{aligned}
\right.
\end{equation}

\begin{lemma}[Solvability of the Riccati equation on a general interval]\label{lm:Riccati}
Let $t_0<1$, $\ell:=1-t_0$, and consider the initial value problem
\[
\frac{\mathrm{d}M}{\mathrm{d}t} + \frac12 M^{2} = N \quad \text{on } [t_0,1],\qquad M(t_0)=0,
\]
with $N\in C^{0}[t_0,1]$.
\begin{itemize}
\item[(i)] If $\|N\|_{C^{0}[t_0,1]} < \dfrac{1}{2\ell^{2}}$, then there exists a unique solution $M\in C^{1}[t_0,1]$ satisfying
\[
\|M\|_{C^{0}[t_0,1]} \le 2\ell\|N\|_{C^{0}[t_0,1]},\qquad
\|M\|_{C^{1}[t_0,1]} \le 2(\ell+1)\|N\|_{C^{0}[t_0,1]}.
\]
\item[(ii)] For any $N_{1},N_{2}$ with $\|N_{i}\|_{C^{0}[t_0,1]} \le \delta < \dfrac{1}{2\ell^{2}}$, let $M_{1},M_{2}$ be the corresponding unique solutions. Then
\[
\|M_{1}-M_{2}\|_{C^{1}[t_0,1]} \le \frac{\ell+1}{1-2\ell^{2}\delta}\,\|N_{1}-N_{2}\|_{C^{0}[t_0,1]}.
\]
\end{itemize}
\end{lemma}

\section{Some discussions}\label{sec:Discussions}

\begin{theorem}\label{thm:A}
Let $\BF$ be an axisymmetric vector field of class $C^1$, homogeneous of degree $-3$ on $\mathbb{R}^3_+$, and let $(\curl\,\BF)^{\tan}$ denote the tangential component of $\curl\,\BF$ on the unit sphere $\mathbb{S}^2$.
Then there exists $\varepsilon>0$ such that if
\[
\|(\curl\,\BF)^{\tan}\|_{C^0(\mathbb{S}^2\cap\mathbb{R}^3_+)} \le \varepsilon,
\]
the system \eqref{eqn:NSEhfD} admits a unique \emph{small axisymmetric self-similar solution in the bounded angular-profile class}. More precisely,
\[
\Bu\in C^{2,\alpha}_{\mathrm{loc}}(\overline{\mathbb R^3_+}\setminus\{0\}),
\qquad p\in C^{1,\alpha}_{\mathrm{loc}}(\overline{\mathbb R^3_+}\setminus\{0\})
\]
for every $\alpha\in(0,1)$.

Moreover, if $\BF$ has no swirl, then any axisymmetric self-similar solution is necessarily swirl-free and unique.
\end{theorem}

\begin{proof}[\textbf{Proof of Theorem~\ref{thm:A} (no-swirl case)}]
The no-swirl case is proved by applying Lemma~\ref{lm:Riccati} with $t_0=0$ and $\ell=1$ to the Riccati equation \eqref{eqn:Riccati}.

\noindent
\underline{Step 1.} \emph{Estimate of $G+2\mcC$.}
For a swirl-free force $\BF = \rho^{-3}(X(\varphi)\Be_\rho + Y(\varphi)\Be_\varphi)$,
\eqref{eqn:sphericalcurl} gives $\nabla\times\BF = -\rho^{-4}(X'+2Y)\Be_\theta$,
and therefore
\[
\|(\curl\BF)^{\tan}\|_{C^0(\mathbb{S}^2\cap\mathbb{R}^3_+)} = \|X'+2Y\|_{C^0[0,\frac{\pi}{2}]}.
\]
Recall that $G(\vp)=X(\vp)+2\int_0^\vp Y(\tau)\, \rmd\tau$ and $\mcC=-\int_0^{\frac{\pi}{2}}G(\tau)\sin\tau \cos\tau \, \rmd\tau$ from \eqref{eqn:constC}.
A direct calculation shows
\begin{align}\label{eqn:Gand2C}
G(\vp)+2\mcC 
= 2\int_0^{\frac{\pi}{2}} \bigl(G(\vp)-G(\tau)\bigr)\sin\tau \cos\tau\, \rmd\tau,
\end{align}
and consequently,
\begin{align*}
\|G+2\mcC\|_{C^0[0,\frac{\pi}{2}]} 
\le \sup_{0 \le \vp,\tau \le \frac{\pi}{2}} |G(\vp)-G(\tau)|
\le \frac{\pi}{2} \|X'+2Y\|_{ C^0[0, \frac{\pi}{2}] }.
\end{align*}
Throughout this proof we keep the coefficients $\pi/2$ explicit because they encode the interval length; this will be convenient in Section~\ref{sec:Discussions} when we replace $[0,\pi/2]$ by $[0,\varphi_0]$ and need to track the angular dependence.

\noindent
\underline{Step 2.} \emph{Estimate of $N$.}
Recall that
\[
N(t)=\frac{-1}{(1-t^2)^2} \int_t^1 J(s) \, \rmd s,
\]
where $J(t)=\int_0^{\vp} (G(\tau) +2\mcC) \sin \tau \, \rmd\tau$ with $t=\cos\vp$.
Using $J(1)=0$, the numerator of $N$ becomes
\begin{equation*}
\begin{aligned}
- \int_t^1 J(s)\, \mathrm{d}s
&= t J(t) + \int_t^1 s \,\frac{\mathrm{d}J}{\mathrm{d}t}(s) \, \mathrm{d}s
= \int_t^1 (s-t)\, \frac{\mathrm{d}J}{\mathrm{d}t}(s)\, \mathrm{d}s \\
&= (1-t)^2 \int_0^1 \frac{\mathrm{d}J}{\mathrm{d}t}\bigl( 1-\theta(1-t) \bigr)(1-\theta) \, \mathrm{d}\theta .
\end{aligned}
\end{equation*}
Therefore
\begin{align}\label{eqn:NandJprime}
N(t) = \frac{1}{(1+t)^2}\int_0^1 \frac{\mathrm{d}J}{\mathrm{d}t}\bigl( 1-\theta(1-t)\bigr)(1-\theta) \, \rmd\theta.
\end{align}
Because $\frac{\mathrm{d}J}{\mathrm{d}t}(t) = -(G(\vp)+2\mcC)$ with $t=\cos\vp$, we have $\bigl\|\frac{\mathrm{d}J}{\mathrm{d}t}\bigr\|_{C^0[0,1]} = \|G+2\mcC\|_{C^0[0,\frac{\pi}2]}$, and
\[
\|N\|_{C^0[0,1]} \le \frac12 \|G+2\mcC\|_{C^0[0,\frac{\pi}{2}]} \le \frac{\pi}{4} \|X'+2Y\|_{ C^0[0, \frac{\pi}{2}] }.
\]

\noindent
\underline{Step 3.} \emph{Existence and uniqueness.}
Set $\varepsilon_1 = \frac{2}{\pi}$. If $\|X'+2Y\|_{C^0[0, \frac{\pi}{2}]} < \varepsilon_1$, then $\|N\|_{C^0[0, 1]} < 1/2$; by Lemma~\ref{lm:Riccati}(i) there exists a unique $M\in C^1[0,1]$. 
Defining $g(\vp)=M(\cos\vp)\sin\vp$, we obtain $g\in C^1[0,\pi/2]$.
To see that $g(\vp)\cot\vp$ is continuous we use $g(0)=0$:
\[
g(\vp)\cot\vp = \cot\vp \int_0^{\vp} g'(\tau)\,\rmd\tau
= \vp\cot\vp \int_0^1 g'(\vp s)\,\rmd s.
\]
Because $|\vp\cot\vp| \le \frac{\pi}{2}$ on $[0,\frac{\pi}{2}]$, the product $g(\varphi)\cot\varphi$ is continuous. Setting $f = -(g' + g\cot\varphi)$ (see \eqref{eqn:halfD2}$_3$) yields $f\in C^0[0,\pi/2]$.  The pressure $P$ is recovered from \eqref{eqn:eqnforP}.  With $h\equiv0$, the quartet $(f,g,h,P)$ satisfies the reduced system \eqref{eqn:halfD1} in the distributional sense. 
Consequently,
\[
\Bu = \frac{f(\varphi)}{\rho}\Be_\rho + \frac{g(\varphi)}{\rho}\Be_\varphi \in C^0(\R^3_+),
\qquad
p = \frac{P(\varphi)}{\rho^2}
\]
fulfil the Navier--Stokes equations \eqref{eqn:NSEhfD} in the sense of distributions.
Uniqueness of $\Bu$ is inherited from the uniqueness of the Riccati equation \eqref{eqn:Riccati}.
Since the angular profiles are $C^1$ and the equations are satisfied away from the origin, standard local Stokes regularity yields \[\n\Bu\in C^{2,\alpha}_{\mathrm{loc}}(\overline{\mathbb R^3_+}\setminus\{0\}),\qquad p\in C^{1,\alpha}_{\mathrm{loc}}(\overline{\mathbb R^3_+}\setminus\{0\})\n\] for every $\alpha\in(0,1)$. Hence all derivatives appearing in the reduction to the ODE system are classical away from the origin.
\end{proof}

\end{document}
